\documentclass[10pt,a4paper]{amsart}
\usepackage[margin=19mm]{geometry}
\usepackage{amsmath,amssymb,mathtools}
\usepackage[scr=rsfs,cal=euler]{mathalfa}
\usepackage{xfrac}
\usepackage[unicode,hidelinks,bookmarks=false]{hyperref}
\newcommand{\ie}{i.e.\ }
\newcommand{\cf}{cf.\ }
\newcommand{\resp}{respectively\ }
\newcommand{\Subsection}{\subsection}
\providecommand{\nobreakdash}{\nobreak\hbox{-}}
\setlength{\emergencystretch}{2em}
\multlinegap0pt
\usepackage{amssymb}
\usepackage{dsfont}
\usepackage[shortlabels]{enumitem}
\usepackage{mathtools}
\usepackage{bm}
\usepackage{xcolor}
\newtheorem{proposition}{Proposition}
\newcommand{\II}{\mathcal{I}}
\newcommand{\N}{\mathds{N}}
\newcommand{\nn}{\nonumber}
\title{SDP bounds on quantum codes: rational certificates}
\author{Gerard Angl\`es Munn\'e, Felix Huber}
\date{Source: arXiv:2603.19901v1 (CC BY 4.0)}
\begin{document}
\maketitle

\noindent\emph{Source and licence.} SDP bounds on quantum codes: rational certificates. Version arXiv:2603.19901v1 (CC BY 4.0), distributed by arXiv under the Creative Commons Attribution 4.0 International licence, http://creativecommons.org/licenses/by/4.0/, as recorded in the arXiv licence metadata for this exact version and shown on the abstract page https://arxiv.org/abs/2603.19901v1. Full licence notice: \texttt{licenses/arxiv-2603.19901-CC-BY-4.0.txt}.\par\smallskip

\noindent\textit{Editorial scope.} Counted unit: the article's proposition rmk:dualpure (source0.tex lines 1054--1071). The article's complete original proof of it is delivered with it (source0.tex lines 1072--1082, one \texttt{proof} environment of 11 lines). No mathematics is written, repaired or re-derived: the statement and the proof are the article's own lines, in the article's own notation, and no retained line is substituted or re-flowed.  Retained uncounted as the source-local context the proof needs: lines 615--630; lines 654--656; lines 1020--1035.  Nothing was omitted from any retained span, so the package carries no omission record.  No retained line contains a citation, so the wrapper carries no bibliography and no bibliography entry.  No retained line contains a figure environment, an image-inclusion command or drawing code, so no image is delivered and the package declares no asset dependency.  Editorial formatting only, in addition: an AMS \texttt{amsart} preamble that restates the article's own declarations and macros used by the retained lines, namely the environments \texttt{proposition} on separate counters, together with the standard packages those lines need.  The retained environments print in this document as Proposition 1 in order of appearance, whereas the article prints the same items with the numbering of its own full document.  The complete native source of the article as distributed in the arXiv e-print for this exact version is bundled unchanged as \texttt{sources/196811/source0.tex}.
\begin{align}\label{eq:sdpx_relax}
	\textnormal{find} 		\quad & x^{t,p}_{i,j}  \nn \\
	\textnormal{subject to} 	\quad
	&  x^{0,0}_{0,0}=1, \nn \\
	%
    &x^{t,p}_{i,j}=0 \quad\quad\quad \text{if}\quad  t-p \quad  \text{is odd}\,,\nn \\
	& x^{t,p}_{i,j}=x^{t',p'}_{i',j'}\,\,\,\, \quad \text{if} \quad  t-p=t'-p' \quad \text{is even} \quad \text{and} \nn \\
	%
	&\quad\quad\quad\quad\quad\quad\,\, (i,j,i+j-t-p) \quad\text{a permutation of} \quad(i',j',i'+j'-t'-p')\,, \nn \\
	%
	& \sum_{\substack{(i,j,t,p)\in \II(n) \\ k=i+j-t-p}}  \gamma^{t,p}_{i,j} x^{t,p}_{i,j}=\frac{2^n}{K} \gamma^{0,0}_{k,0} x^{0,0}_{k,0}   \,, \nn\\
	%
	& K2^{-n}\sum^n_{i=0} K_j(i,n) \gamma^{0,0}_{i,0} x^{0,0}_{i,0}\quad = \quad \gamma^{0,0}_{j,0} x^{0,0}_{j,0} \quad\quad \text{for} \quad 0<j<\delta \,, \nn \\
	%
	&\bigoplus_{\substack{a,k \in \N_0\\ 0\leq a\leq k\leq n+a-k}} \left(\sum\limits_{\substack{t,p\in \N_0 \\ 0 \leq p \leq t \leq i,j \\ i+j\leq t+n}}\alpha(i,j,t,p,a,k)x_{i,j}^{t,p}\right)_{i,j=k}^{n+a-k}   \succeq 0\,,
\end{align}

\begin{align}\label{eq:klred_pure_xijtp}
x^{t,p}_{i,j}=0 \quad\quad \text{if} \quad\quad \{i,j,i+j-t-p\} \cap \{1,\cdots, \delta-1\} \neq \emptyset\,.
\end{align}

\begin{align}\label{eq:dual_solall}
	\alpha \quad = \quad \max_{Y^{(a,k)},\, C_i} \quad &   (KD_0-C_0)  - {y^{0,0}_{0,0}} \quad &  \nn \\
	\text{subject to}  \quad & Y^{(a,k)} \succeq0,\nn \\
	& y^{t,p}_{i,j}\ = \frac{1}{\gamma^{t,p}_{i,j}} \sum^{\min(i,j)}_{k=0} \sum^k_{a=\max(i,j)+k-n} \alpha(i,j,t,p,a,k) Y^{(a,k)}_{i-k,j-k}\,, \nn \\
	&D_i = 2^{-n} \sum^{\delta-1}_{j=0} K_j(i,n)  C_{j} \,, \nn \\
	%
	&Q_i=\frac{1}{\left(\frac{2^n}{K}-2\right)}   (2y^{0,0}_{i,0} + y^{i,i}_{i,i} -KD_i + C_i)
	\quad\quad\quad\quad\quad\hspace{-0.03cm} \text{for} \quad 0 <i <\delta \,,\nn \\
	%
	& Q_i=\frac{1}{\left(\frac{2^n}{K}-2\right)}   (2y^{0,0}_{i,0} + y^{i,i}_{i,i} -KD_i)
	\quad\quad\quad\quad\quad\quad\quad\, \text{for} \quad \delta \leq i \leq n \,, \nn \\
	& \sum_{\substack{ (i',j',t',p')\in \II(4,n) \\ (t-p)=(t'-p') \\ (i',j',i'+j'-t'-p') \, \text{a} \\
			\, \text{permutation of} \,  (i,j,i+j-t-p) \\ }} \!\!\!\!\!\!\!\!\!\!  
	\left(y^{t',p'}_{i',j'}+ Q_{i'+j'-t'-p'}\right) = 0 \quad  \hspace{2.3em} \text{if} \quad (t-p) \quad \text{is even}  \nn \\ &  \vspace{-10cm}\hspace{6cm}
    \quad\quad\quad\quad \text{and} \quad i,j,i+j-t-p \neq 0 \,.
\end{align}

\begin{proposition}\label{rmk:dualpure}
    For pure codes 
    the dual of SDP 
~\eqref{eq:sdpx_relax} 
with the additional constraint
Eq.~\eqref{eq:klred_pure_xijtp} 
consists of SDP~\eqref{eq:dual_solall}
with the following modifications:

% 
\begin{itemize}
    \item[(i)] The variable $Q_i$ is indexed by $\delta \leq i \leq n$ 
    instead of $0 \leq i \leq n$.
    \item[(ii)] Last constraint in SDP~\eqref{eq:dual_solall} 
    is only applied when $(i,j,t,p)$ satisfies that $(t-p)$ is even and $i\geq \delta$ and $j\geq \delta$ and $i+j-t-p\geq \delta$.
\end{itemize}

\end{proposition}

\begin{proof}
For pure codes the conditions of Eq.~\eqref{eq:klred_pure_xijtp} can be equivalently written as: $x^{t,p}_{i,j} = 0$ if $0<i<\delta$ or $0<j<\delta$ or $0<i+j-t-p<\delta$. 
% 
By SDP duality, each primal variable corresponds to a dual constraint.
% 
Thus, all constraints from the dual corresponding to the removed primal variables can be omitted. 
% 
In our case, (i) the last equality of Eq.~\eqref{eq:dual_solall} is only applied if $(t-p)$ is even, $i\geq \delta$, $j\geq \delta$, and $i+j-t-p\geq \delta$.
% 
Furthermore, (ii) $Q_i$ is now only defined when $\delta \leq i \leq n$.  This ends the proof.
\end{proof}

\end{document}
